\paragraph{Bind a scheduled decision when its date arrives.}
The preceding sequence fixes the state read by each decision, but its caller
must supply the corresponding future book hashes. An authored simulation should
instead supply its starting book and a calendar, then calculate those later
states as it proceeds. The bounded \texttt{StateDecisionSchedule/1} extension
does this by materializing each decision request only at its turn. Its expected
state hash is the hash of the current book, after earlier successful events.

A schedule supplies a named rule, explicit priority and either a list of
calendar dates or a start date, exclusive end date and positive integer number
of days between occurrences. For example, daily decisions from 2 January 2027
up to, but excluding, 5 January produce the same three decisions as the explicit
dates 2, 3 and 4 January. Calendar arithmetic includes leap days. It introduces
no month-end adjustment, business calendar, timezone conversion or contractual
day-count convention. All occurrences must fall within the current fiscal
period and no earlier than the book's last event.

Different decisions on the same day need distinct explicit priorities. An
unresolved tie stops scheduling before economic execution; rearranging the input
array does not choose an order. Each decision and sale receives a stable identity
derived from the declared root and schedule names and its occurrence number.
The complete schedule also retains its own content identity, so an equivalent
list of dates and recurrence can produce identical execution while remaining
different authored documents. Every declared policy is statically checked
before the first event, including a later policy that an earlier refusal might
otherwise prevent the runner from reaching.

The checked calendar reproduces the manually assembled stock/reserve sequence:
one unit sold, then half a unit, then no action. Cash is USD~474.50. The selected
lot retains USD~5 of carrying amount, while the inventory control totals USD~65:
the opening USD~80 also included USD~60 outside that lot. With a zero reserve,
the same rule sells both units over two decisions and cash reaches USD~482;
the next decision observes depletion and takes no action. Distinguishing a
selected lot from the company's inventory control is necessary when checking
the full balance sheet.

An unresolved or refused decision stops the sequence with the exact current
book, earlier outcomes and pending calendar retained. The stop record includes
the materialized request, its state hash and the proposed event where available.
A no-action advances the decision cursor without changing the book's event
date. The scheduler accepts up to 32 schedules and 64 total decisions; the
existing request checker also bounds serialized size, nesting and structure.
These limits keep the reference finite and do not establish production capacity.

Eighteen focused scheduler methods, 28 direct decision regression methods and
49 exact comparisons across five cases preserve manually assembled requests,
complete native outcomes and all financial report fields. The evidence is in
\path{data/accounting-master/state-decision-schedule-20261025/}. The calendar
retains the explicit premise that no external economic event is omitted.
Combining decisions with separately supplied financing, tax, purchase or period
events requires those events' own timing and accounting rules. General calendar
generation, historical completeness and durable restart remain separate work.

